% Generated draft: requires mathematical and source review.
\documentclass[11pt]{article}
\usepackage[margin=25mm]{geometry}
\usepackage{fontspec}
\setmainfont{Latin Modern Roman}
\usepackage{amsmath,amssymb,mathtools,unicode-math}
\setmathfont{Latin Modern Math}
\usepackage{enumitem,needspace,xurl,hyperref}
\setlength{\parindent}{0pt}
\setlength{\parskip}{5pt}
\setlength{\emergencystretch}{3em}
\newcommand{\N}{\mathbb{N}}
\newcommand{\Z}{\mathbb{Z}}
\newcommand{\Q}{\mathbb{Q}}
\newcommand{\R}{\mathbb{R}}
\newcommand{\C}{\mathbb{C}}
\newcommand{\F}{\mathbb{F}}
\newcommand{\A}{\mathbb{A}}
\newcommand{\PP}{\mathbb{P}}
\newcommand{\E}{\mathbb{E}}
\newcommand{\eps}{\varepsilon}
\newcommand{\dd}{\,\mathrm{d}}
\begin{document}
\section*{UnsolvedMath ID: 30006164}
\textbf{Source problem code:} catalogue entry
\subsection{Definitions and mathematical statement}
Assume \(\mathrm{ZF}+\mathrm{AD}+V=L(\mathbb R)\). Here \(\mathrm{AD}\) asserts determinacy for every perfect-information game of length \(\omega\) in which the players play natural numbers, and \(L(\mathbb R)\) is the constructible universe over the reals. Write \(J_\gamma(\mathbb R)\) for the Jensen hierarchy over the reals, with the indexing and fine-structural conventions of the cited report.

Consider any ordinals satisfying either of the following alternatives:
\begin{itemize}
\item \([\alpha,\gamma]\) is a weak \(\Sigma_1\) gap;
\item \([\alpha,\beta]\) is a strong gap and \(\gamma=\beta+1\).
\end{itemize}
Let \(e<\omega\) be least such that
\[
\rho_{e+1}^{J_\gamma(\mathbb R)}=\mathbb R.
\]
The formulation presupposes that such an \(e\) exists. The projectum notation and the weak and strong gap terminology are used in the report's sense; their precise definitions are not supplied in the record.

For this draft, a \emph{mouse set} means a set of reals of the form \(\mathbb R\cap M\), where \(M\) is a mouse in the intended sense of the report. The required mouse conventions, including iterability requirements, are not specified in the supplied excerpt. Likewise, \(\mathrm{OD}^{*\mathbb R}_{\gamma,e+2}\) denotes the source-designated collection of reals: the record does not define its definability, parameter, or starred-notation conventions.

\emph{Question as given in the reconstructed statement.} For every choice satisfying the preceding hypotheses, does there exist a mouse \(M\), in the report's intended sense, such that
\[
\mathrm{OD}^{*\mathbb R}_{\gamma,e+2}=\mathbb R\cap M?
\]
The assertion concerns the entire displayed collection, without restricting to a cone of reals. It does not merely ask for the parameterized assertion on a cone.

\emph{Unresolved source qualification.} The reconstructed statement uses \(\mathrm{OD}^{*\mathbb R}_{\gamma,e+2}\), whereas the original excerpt states the open question using an unstarred \(\mathrm{OD}_{\gamma,e+2}\). The supplied information does not establish their equality at this level. Thus the displayed question preserves the reconstructed target but cannot resolve whether it is exactly the target of the original excerpt. A fully specified mathematical statement requires checking these definitions and this discrepancy against the report.
\subsection{Short English statement}
Under \(\mathrm{ZF}+\mathrm{AD}+V=L(\mathbb R)\), the reconstructed question asks whether the indicated collection of definable reals at level \(e+2\) is exactly the set of reals belonging to some mouse, for every permitted gap configuration, without a cone restriction. The starred versus unstarred target and several technical conventions remain unresolved. This is a statement-only draft requiring human mathematical review.
\subsection{Sources}
\begin{itemize}
\item Local UnsolvedMath registry record: permanent ID \(30006164\), source code \texttt{catalogue entry}, title \emph{Mouse-Set Property Without a Cone Restriction}; supplied from the project's problem-registry corpus, \texttt{data/unsolvedmath/problems.json} and supplementary \texttt{catalog/problems.json}. Registry membership in either particular file has not been independently checked here.
\item Supplied report citation: \emph{Set Theory} (2025), Oberwolfach Reports, pp.~75--140, DOI \texttt{10.4171/owr/2025/2}. \url{https://doi.org/10.4171/owr/2025/2}. The record attributes the problem to the workshop-report collection; no individual proposer is supplied.
\item Supplied literature source, \emph{Oberwolfach Report 2/2025}: \url{https://content.ems.press/assets/public/full-texts/serials/owr/22/1/14299082/online/10.4171-owr-2025-2.pdf}.
\item Supplied literature source: Schlutzenberg, \emph{The Mouse Set Theorem Just Past Projective}. \url{https://arxiv.org/abs/2307.08856}.
\item Literature assessment dated exactly \texttt{2026-08-22}: the supplied curation labels the problem open and reports that 2024 mouse-set results settle higher complexity levels and cone versions, while the unrestricted case at the extracted level remains open. This is supplied dated curation, not an independently verified literature assessment. The record itself requests independent verification of time-sensitive status and cited claims.
\item The registry labels the reconstructed statement \texttt{corrected\_verified} and says it was checked against the official report. That verification claim is supplied metadata, not a verification performed for this draft. The excerpt also cites Rudominer and Steel as reference [1] and a 2024 theorem attributed to ``S.'' as reference [4]; complete bibliographic entries for those numbered references are not supplied.
\end{itemize}
\end{document}
